\chapter{Posta elettronica e applicazioni peer-to-peer}

% ══════════════════════════════════════════════════════════════
\section{La posta elettronica}
% ══════════════════════════════════════════════════════════════

La posta elettronica (\emph{e-mail}) è una delle applicazioni più antiche e importanti di Internet.
È realizzata in modo \textbf{client-server}, con due tipi di host:

\begin{itemize}
    \item lo \textbf{user agent}: l'applicazione con cui l'utente gestisce la posta (Thunderbird, K-9 Mail,
          l'app Mail del telefono, \dots);
    \item il \textbf{mail server}: un server che conserva la posta e mantiene una \textbf{casella}
          (\emph{mailbox}) per ciascun utente.
\end{itemize}

A differenza di HTTP, qui la comunicazione avviene \textbf{anche tra server}. Servono quindi due
applicazioni (e due protocolli): una tra user agent e mail server, e una tra mail server e mail server.

\begin{figure}[H]
\centering
\begin{tikzpicture}[every node/.style={font=\footnotesize\sffamily},
  ms/.style={draw=netblue, fill=cloudbg, rounded corners, minimum width=2.4cm, minimum height=1.5cm, align=center}]
  \node[ms] (s1) at (0,0) {\icon[0.5cm]{server}\\mail server\\di Alice};
  \node[ms] (s2) at (7,0) {\icon[0.5cm]{server}\\mail server\\di Bob};
  \node[ms] (s3) at (3.5,-2.6) {\icon[0.5cm]{server}\\altro\\mail server};
  \node (ua1) at (-3.4,1.2) {\icon[1cm]{laptop}}; \node[below=-2pt of ua1] {user agent (Alice)};
  \node (ua2) at (10.4,1.2) {\icon[0.7cm]{smartphone}}; \node[below=-2pt of ua2] {user agent (Bob)};
  \node (ua3) at (-3.4,-2.2) {\icon[0.9cm]{pc}}; \node[below=-2pt of ua3] {user agent};
  \draw[msg, netred, very thick] (ua1) -- node[above, sloped] {SMTP} (s1);
  \draw[msg, netred, very thick] (s1) -- node[above] {SMTP} (s2);
  \draw[msg, netgreen, very thick] (s2) -- node[above, sloped] {POP3 / IMAP / HTTP} (ua2);
  \draw[msg, netred, {Stealth}-{Stealth}] (s1) -- node[above, sloped] {SMTP} (s3);
  \draw[msg, netred, {Stealth}-{Stealth}] (s2) -- node[above, sloped] {SMTP} (s3);
  \draw[{Stealth}-{Stealth}, thick, netgreen] (ua3) -- (s3);
  \foreach \s in {s1, s2, s3} \draw[fill=llink] ([xshift=-0.35cm, yshift=-0.55cm]\s.south east) rectangle ++(0.3,0.25);
  \node[lbl] at (8.9,-1.2) {\colorbox{llink}{\strut}\ casella di Bob};
\end{tikzpicture}
\caption{Il sistema di posta: SMTP per inviare (dagli user agent ai server e tra server), protocolli di accesso per leggere la posta dalla propria casella.}
\end{figure}

Gli utenti non si scambiano la posta direttamente (come nella messaggistica istantanea): si
preferisce passare per i mail server, più affidabili e specializzati.

% ══════════════════════════════════════════════════════════════
\section{SMTP}
% ══════════════════════════════════════════════════════════════

\dfn{SMTP}{
    Il \textbf{Simple Mail Transfer Protocol} (SMTP) è il principale protocollo applicativo usato per
    \textbf{trasferire la posta tra mail server}. Ha un lato client (chi invia) e un lato server (chi
    riceve), entrambi eseguiti sui mail server, e usa il trasporto affidabile di \textbf{TCP} sulla
    \textbf{porta 25}.
}

Quando un nuovo messaggio arriva sul server, viene messo nella casella del destinatario, in attesa
che lo user agent lo scarichi.

\subsection{Il flusso completo di una mail}

Alice invia una mail a Bob (\texttt{bob@someschool.edu}). Sono coinvolti quattro elementi: lo user agent
di Alice, il mail server di Alice, il mail server di Bob, lo user agent di Bob.

\begin{figure}[H]
\centering
\begin{tikzpicture}[every node/.style={font=\footnotesize\sffamily}]
  \node (a) at (0,0) {\icon[1cm]{laptop}}; \node[below=-2pt of a] {user agent di Alice};
  \node[draw=netblue, fill=cloudbg, rounded corners, minimum width=2.6cm, minimum height=1.6cm] (sa) at (4,0) {};
  \node at (4,0.45) {mail server di Alice};
  \node[draw, fill=llink, minimum width=1.4cm, minimum height=0.4cm] at (4,-0.3) {coda messaggi};
  \node[draw=netblue, fill=cloudbg, rounded corners, minimum width=2.6cm, minimum height=1.6cm] (sb) at (9.5,0) {};
  \node at (9.5,0.45) {mail server di Bob};
  \node[draw, fill=llink, minimum width=1.4cm, minimum height=0.4cm] at (9.5,-0.3) {casella di Bob};
  \node (b) at (13.5,0) {\icon[0.7cm]{smartphone}}; \node[below=-2pt of b] {user agent di Bob};
  \draw[msg, very thick] (a) -- node[above] {\circled{2}} (sa);
  \draw[msg, very thick, netred] (sa) -- node[above] {\circled{3}\ \circled{4}\ SMTP su TCP} (sb);
  \draw[msg, very thick, netgreen] (sb) -- node[above] {\circled{6}} (b);
  \node[above=6pt of a] {\circled{1}};
  \node[above=-2pt of sb.north] {\circled{5}};
\end{tikzpicture}
\caption{I sei passi dell'invio di una mail da Alice a Bob.}
\end{figure}

\begin{enumerate}
    \item Alice apre il suo user agent, inserisce l'indirizzo di Bob, scrive il messaggio e chiede di inviarlo.
    \item Lo user agent di Alice invia il messaggio al \textbf{suo} mail server, dove viene messo in una
          \textbf{coda di messaggi}.
    \item Il lato client di SMTP, sul server di Alice, vede il messaggio in coda e \textbf{apre una connessione
          TCP} verso il server SMTP in esecuzione sul mail server di Bob.
    \item Dopo un handshaking SMTP iniziale, il client SMTP invia il messaggio di Alice sulla connessione TCP.
    \item Sul server di Bob il lato server di SMTP riceve il messaggio e lo mette nella \textbf{casella di Bob}.
    \item Bob apre il proprio user agent e legge il messaggio quando vuole.
\end{enumerate}

\info{Nessun server intermedio}{
    SMTP di norma \textbf{non} usa mail server intermedi: la connessione TCP va direttamente dal server di
    Alice a quello di Bob, anche se sono ai lati opposti del pianeta. Se il server di Bob è spento, il
    messaggio resta nella coda del server di Alice, che riprova più tardi.
}

\subsection{Dettagli del protocollo}

\begin{itemize}
    \item Il client SMTP (sul server mittente) apre una connessione TCP verso la \textbf{porta 25} del server
          SMTP (sul server destinatario).
    \item Se il server non risponde, il client \textbf{riprova più tardi}.
    \item Stabilita la connessione, client e server fanno un \textbf{handshaking applicativo}: si presentano
          prima di trasferire informazioni. In questa fase il client indica l'indirizzo del \textbf{mittente} e
          quello del \textbf{destinatario}.
    \item Poi il client invia il messaggio, contando sul trasferimento affidabile di TCP per farlo arrivare
          senza errori.
\end{itemize}

\begin{esempio}{Un dialogo SMTP}
Le righe che iniziano con \texttt{S:} sono inviate dal server, quelle con \texttt{C:} dal client
(da Kurose). Come HTTP, anche SMTP usa comandi testuali leggibili.
\begin{lstlisting}
S: 220 hamburger.edu
C: HELO crepes.fr
S: 250 Hello crepes.fr, pleased to meet you
C: MAIL FROM: <alice@crepes.fr>
S: 250 alice@crepes.fr ... Sender ok
C: RCPT TO: <bob@hamburger.edu>
S: 250 bob@hamburger.edu ... Recipient ok
C: DATA
S: 354 Enter mail, end with "." on a line by itself
C: Do you like ketchup?
C: How about pickles?
C: .
S: 250 Message accepted for delivery
C: QUIT
S: 221 hamburger.edu closing connection
\end{lstlisting}
\end{esempio}

% ══════════════════════════════════════════════════════════════
\section{Accedere alla posta}
% ══════════════════════════════════════════════════════════════

Usare dei server SMTP è più affidabile della consegna diretta da utente a utente:

\begin{itemize}
    \item i server sono, per definizione, \textbf{sempre accesi} e visibili a tutti i dispositivi della rete
          (il portatile di Bob, invece, può essere spento);
    \item i server sono \textbf{certificati};
    \item il server può \textbf{riprovare continuamente} a consegnare la posta se la consegna fallisce, un
          processo piuttosto oneroso da fare sul proprio dispositivo.
\end{itemize}

D'altra parte, per \textbf{leggere} la posta dal proprio server serve un'altra applicazione
client-server, con un protocollo diverso da SMTP (SMTP è un protocollo ``push'': serve a spingere la
posta verso un server, non a tirarla giù). I più diffusi sono POP3, IMAP e HTTP.

\subsection{POP3}

\dfn{POP3}{
    Il \textbf{Post Office Protocol --- Version 3} (POP3) è un protocollo di accesso alla posta
    estremamente semplice. Lo user agent (client) apre una connessione TCP verso il mail server sulla
    \textbf{porta 110}.
}

Stabilita la connessione, POP3 attraversa tre fasi:

\begin{figure}[H]
\centering
\begin{tikzpicture}[every node/.style={font=\footnotesize\sffamily}, ph/.style={draw, rounded corners, fill=#1, minimum width=3.6cm, minimum height=1.6cm, align=center, text width=3.4cm}]
  \node[ph=lapp!60] (a) at (0,0) {\textbf{1. Autorizzazione}\\[2pt] \scriptsize invio di nome utente e password};
  \node[ph=ltra!60] (t) at (4.8,0) {\textbf{2. Transazione}\\[2pt] \scriptsize scaricare messaggi, marcarli per la cancellazione, statistiche};
  \node[ph=lnet!60] (u) at (9.6,0) {\textbf{3. Aggiornamento}\\[2pt] \scriptsize dopo \texttt{quit}: il server cancella i messaggi marcati};
  \draw[msg, very thick] (a) -- (t); \draw[msg, very thick] (t) -- (u);
\end{tikzpicture}
\caption{Le tre fasi di una sessione POP3.}
\end{figure}

\begin{enumerate}
    \item \textbf{Autorizzazione}: lo user agent invia nome utente e password per autenticare l'utente.
    \item \textbf{Transazione}: lo user agent può recuperare messaggi, marcarli per la cancellazione, togliere
          le marcature e ottenere statistiche sulla casella.
    \item \textbf{Aggiornamento}: quando il client dà il comando \texttt{quit}, la sessione termina e il server
          cancella i messaggi marcati.
\end{enumerate}

\subsection{IMAP}

Con POP3 i messaggi si possono solo scaricare o cancellare: cercare, organizzare in cartelle e simili va
fatto sulla macchina locale. Se si legge la posta da più dispositivi, questo è un problema.

\dfn{IMAP}{
    L'\textbf{Internet Mail Access Protocol} (IMAP, porta 143) permette al server di offrire funzioni
    aggiuntive, mantenendo la posta \textbf{sul server}:
    \begin{itemize}
      \item creare e gestire \textbf{cartelle} sul server;
      \item \textbf{cercare} nelle cartelle remote i messaggi che soddisfano certi criteri;
      \item scaricare solo \textbf{parti} di un messaggio (solo le intestazioni, solo gli allegati, \dots);
      \item collegare \textbf{più client} contemporaneamente alla stessa casella.
    \end{itemize}
}

\subsection{HTTP (webmail)}

Sempre più utenti inviano e leggono la posta dal \textbf{browser}. L'accesso via Web fu introdotto da
Hotmail a metà degli anni '90 ed è oggi l'approccio più diffuso (Google, Yahoo!, Libero, e quasi tutte le
università e le aziende, UniNA compresa). In questo caso lo user agent è un normale browser, e l'utente
comunica con la propria casella remota via \textbf{HTTP} invece che con POP3 o IMAP.

\warning{Tra server resta SMTP}{
    HTTP sostituisce POP3/IMAP solo tra \textbf{user agent e server}. Tra \textbf{mail server} la posta
    continua a viaggiare con SMTP.
}

\begin{table}[H]
\centering
\begin{tabular}{llll}
\toprule
\textbf{Protocollo} & \textbf{Tra} & \textbf{Porta TCP} & \textbf{Caratteristica} \\
\midrule
SMTP & server $\leftrightarrow$ server, user agent $\rightarrow$ server & 25 & invio (push) \\
POP3 & user agent $\leftarrow$ server & 110 & scarica e (di solito) cancella \\
IMAP & user agent $\leftrightarrow$ server & 143 & la posta resta sul server, cartelle, ricerca \\
HTTP & browser $\leftrightarrow$ server & 80/443 & webmail \\
\bottomrule
\end{tabular}
\caption{I protocolli della posta elettronica.}
\end{table}

% ══════════════════════════════════════════════════════════════
\section{Peer-to-peer}
% ══════════════════════════════════════════════════════════════

A differenza del client-server, un'architettura P2P usa i server in modo minimo, o per nulla: coppie di
host collegati in modo intermittente (i \textbf{peer}) comunicano direttamente. I peer non appartengono
a un fornitore di servizi: sono computer e portatili controllati dagli utenti.

\subsection{Distribuire un file: due modi}

Un file, un server, e molti host che ne vogliono una copia.

\begin{figure}[H]
\centering
\begin{tikzpicture}[every node/.style={font=\footnotesize\sffamily}]
  \begin{scope}
    \node[font=\small\bfseries\sffamily] at (0,2.4) {Client-server};
    \node (srv) at (0,1.4) {\icon[0.7cm]{server}};
    \node[right=1pt of srv, text=netred] {server: carica a $u_s$};
    \foreach \x [count=\i] in {-2.4,-1.2,0,1.2,2.4} {
      \node (c\i) at (\x,-0.8) {\icon[0.65cm]{laptop}};
      \draw[-{Stealth}, thick, netred] (srv) -- node[fill=white, inner sep=0.5pt, font=\scriptsize\sffamily] {$F$} (c\i);
    }
    \node[lbl, text width=5.2cm] at (0,-2.0) {il server invia \textbf{$N$ copie} complete: $N F$ bit, tutti con la sua sola banda};
  \end{scope}
  \begin{scope}[xshift=8cm]
    \node[font=\small\bfseries\sffamily] at (0,2.4) {Peer-to-peer};
    \node (srv) at (0,1.4) {\icon[0.7cm]{server}};
    \node[right=1pt of srv, text=netred] {server};
    \foreach \x/\y [count=\i] in {-2.7/-1.0, -1.5/0.35, 0/-0.45, 1.5/0.35, 2.7/-1.0} {
      \node (p\i) at (\x,\y) {\icon[0.6cm]{laptop}};
    }
    \draw[-{Stealth}, thick, netred, shorten >=2pt] (srv) -- node[right, font=\scriptsize\sffamily] {$\approx F$} (p3);
    \foreach \a/\b in {3/2, 3/4, 2/1, 4/5} \draw[-{Stealth}, thick, netgreen!60!black, shorten >=3pt, shorten <=3pt] (p\a) -- (p\b);
    \draw[-{Stealth}, thick, netgreen!60!black, shorten >=3pt, shorten <=3pt] (p1) to[bend right=25] (p3);
    \draw[-{Stealth}, thick, netgreen!60!black, shorten >=3pt, shorten <=3pt] (p5) to[bend left=25] (p3);
    \node[lbl, text width=5.2cm] at (0,-2.0) {il server invia \textbf{circa una copia}; i peer si passano i pezzi, ciascuno con la propria banda di upload $u_i$};
  \end{scope}
\end{tikzpicture}
\caption{A sinistra il server invia $N$ copie; a destra ne invia circa una, e i peer si scambiano i pezzi tra loro.}
\end{figure}

Il ragionamento si può rendere quantitativo. Un file di $F$ bit va distribuito a $N$ peer; il server carica a
velocità $u_s$, il peer $i$ carica a $u_i$, e $d_{\min}$ è la velocità di download del peer più lento.

\dfn{Tempo minimo di distribuzione}{
    $$D_{\text{cs}} = \max\left\{\frac{N F}{u_s},\ \frac{F}{d_{\min}}\right\} \qquad\qquad
      D_{\text{P2P}} = \max\left\{\frac{F}{u_s},\ \frac{F}{d_{\min}},\ \frac{N F}{u_s + \sum_{i=1}^{N} u_i}\right\}$$
}

\begin{itemize}
    \item \textbf{Client-server}: il server da solo deve inviare $N F$ bit alla velocità $u_s$, e nessun client può
          finire prima di aver scaricato $F$ bit alla propria velocità (il più lento a $d_{\min}$). Il tempo
          \textbf{cresce linearmente} con $N$, senza limite.
    \item \textbf{P2P}: il server deve inviare almeno una copia ($F/u_s$); il peer più lento non può scaricare più
          in fretta di $F/d_{\min}$; e i $N F$ bit complessivi vengono caricati usando la banda del server
          \textbf{più quella di tutti i peer}. Ogni nuovo peer aggiunge un $F$ al numeratore, ma anche il suo $u_i$
          al denominatore.
\end{itemize}

\begin{figure}[H]
\centering
\begin{tikzpicture}
  \begin{axis}[width=9cm, height=5.4cm, xlabel={numero di peer $N$}, ylabel={tempo minimo (ore)},
      xmin=0, xmax=35, ymin=0, ymax=3.6, legend pos=north west, legend style={font=\scriptsize\sffamily},
      tick label style={font=\scriptsize}, label style={font=\footnotesize\sffamily}, grid=major, grid style={gray!20}]
    \addplot[very thick, netred, domain=0:35, samples=50] {x/10};
    \addlegendentry{client-server}
    \addplot[very thick, netblue, domain=0:35, samples=80] {max(0.1, x/(10+x))};
    \addlegendentry{P2P}
  \end{axis}
\end{tikzpicture}
\caption{Tempo minimo di distribuzione con $u_s = 10u$ (tutti i peer caricano a $u$), $F/u = 1$ ora e $d_{\min}$ abbondante (da Kurose): il client-server cresce come $N/10$, il P2P come $N/(10+N)$ e non supera mai 1 ora.}
\end{figure}

Mille volte i peer sono mille volte l'attesa nel client-server, e quasi nessuna attesa in più nel P2P: è la
\textbf{auto-scalabilità} del P2P, e un esercizio numerico completo si trova nella parte degli esercizi.

\info{Pro e contro del P2P}{
    Nella condivisione di file il P2P \textbf{scala} di solito molto meglio del client-server. D'altra parte è
    più \textbf{complesso da implementare}, e avere tanti client in comunicazione diretta tra loro pone
    problemi di \textbf{sicurezza}.
}

% ══════════════════════════════════════════════════════════════
\section{BitTorrent}
% ══════════════════════════════════════════════════════════════

Un protocollo P2P molto diffuso per la distribuzione di file è \textbf{BitTorrent} (circa 150--170
milioni di utenti stimati nel 2023).

\dfn{Torrent e chunk}{
    Un \textbf{torrent} è l'insieme di tutti i peer che partecipano alla distribuzione di un certo file.
    I peer di un torrent si scambiano \textbf{chunk}, pezzi del file di uguale dimensione
    (tipicamente 256~kB).
}

\subsection{Il ciclo di vita di un peer}

Quando un peer entra in un torrent, senza avere ancora alcun chunk:

\begin{itemize}
    \item comincia ad \textbf{accumulare chunk};
    \item mentre scarica, \textbf{carica} anche chunk verso altri peer;
    \item quando ha l'intero file può andarsene (egoisticamente) oppure restare (altruisticamente) a
          continuare a caricare chunk per gli altri.
\end{itemize}

I peer possono lasciare il torrent in qualsiasi momento, anche con solo una parte dei chunk, e
rientrare più tardi.

\subsection{Il tracker}

\dfn{Tracker}{
    Ogni torrent ha un nodo di infrastruttura, il \textbf{tracker}, che tiene traccia dei peer che
    partecipano al torrent. È, in sostanza, un \textbf{server}, in un'architettura che avrebbe dovuto farne a meno:
    un altro esempio di architettura ibrida.
}

\begin{figure}[H]
\centering
\begin{tikzpicture}[every node/.style={font=\footnotesize\sffamily}]
  \node (tr) at (0,2.8) {\icon[0.7cm]{server}}; \node[right=2pt of tr, text=netred] {tracker};
  \node[ellipse, minimum width=9.5cm, minimum height=3.8cm, fill=netgreen!8, draw=netgreen!60, thick] at (0.4,-0.6) {};
  \node[text=netgreen!50!black, font=\small\bfseries\sffamily] at (-3.2,0.9) {torrent};
  \node (alice) at (-6.2,0) {\icon[1cm]{laptop}}; \node[below=-2pt of alice, text=netblue] {nuovo peer (Alice)};
  \foreach \x/\y/\ic [count=\i] in {-2.2/0.3/pc, 0/0.6/laptop, 2.4/0.2/pc, -1.2/-1.6/laptop, 1.3/-1.7/pc, 3.3/-1.1/laptop} {
    \node (p\i) at (\x,\y) {\icon[0.6cm]{\ic}};
  }
  \draw[msg, dashed, netred] (alice) to[bend left=20] node[above, sloped] {1. registrazione} (tr);
  \draw[msg, dashed, netred] (tr) to[bend left=10] node[below, sloped] {2. lista di peer} (alice);
  \foreach \i in {1,4,2} \draw[{Stealth}-{Stealth}, thick, netblue] (alice) -- (p\i);
  \node[text=netblue] at (-4.2,-1.8) {3. connessioni TCP};
  \foreach \a/\b in {1/2, 2/3, 3/6, 4/5, 5/6, 1/4, 2/5} \draw[thick, black!35] (p\a) -- (p\b);
  \foreach \i/\c in {1/lapp, 2/ltra, 3/lnet} \fill[\c] (5.2+\i*0.35,2.4) rectangle ++(0.3,0.3);
  \node[right] at (6.4,2.55) {chunk del file};
\end{tikzpicture}
\caption{Un nuovo peer si registra presso il tracker, riceve un sottoinsieme di peer e apre con loro connessioni TCP.}
\end{figure}

Quando un nuovo peer entra in un torrent:

\begin{itemize}
    \item si \textbf{registra} presso il tracker e lo informa periodicamente del proprio stato;
    \item il tracker gli fornisce gli indirizzi IP di un \textbf{sottoinsieme casuale} di peer del torrent;
    \item avuta la lista, il nuovo peer prova ad aprire \textbf{connessioni TCP} concorrenti con tutti i peer
          della lista, che diventano i suoi \textbf{vicini} (\emph{neighbors});
    \item durante l'esecuzione alcuni vicini se ne vanno, mentre altri peer (fuori dalla lista iniziale)
          possono chiedere di connettersi.
\end{itemize}

In ogni istante ogni peer ha un sottoinsieme diverso dei chunk. Periodicamente un peer chiede a ogni
vicino (sulle connessioni TCP) la lista dei chunk che possiede.

\subsection{Le politiche del client}

\dfn{Rarest first}{
    Chi \textbf{scarica} decide quali chunk chiedere, e a chi, secondo il principio \textbf{rarest first}:
    si danno priorità ai chunk che hanno \textbf{meno copie} tra i vicini. Così i chunk rari vengono
    ridistribuiti più rapidamente e il numero di copie di ciascun chunk nel torrent tende a equilibrarsi.
}

Chi \textbf{carica} decide a quali richieste rispondere secondo due principi intrecciati:

\begin{itemize}
    \item \textbf{Scambio} (\emph{trading}, detto anche \emph{tit-for-tat}): si dà priorità ai
          \textbf{4 vicini migliori}, cioè quelli che in quel momento ci stanno inviando dati alla velocità più
          alta. Il controllo viene ripetuto ogni \textbf{10 secondi} e la lista dei 4 migliori aggiornata.
    \item \textbf{Selezione casuale} (\emph{optimistic unchoking}): ogni \textbf{30 secondi} si sceglie anche un
          vicino \textbf{a caso} e gli si inviano chunk. Se lo scambio va bene (i due sono buoni partner), i due
          entrano nelle rispettive liste dei migliori. Questo permette anche ai peer con velocità di upload
          compatibili di trovarsi, e ai nuovi arrivati di ricevere i primi chunk.
\end{itemize}

\begin{esempio}{Perché la selezione casuale}
Alice ha appena iniziato e non ha nulla da offrire: nessuno la metterebbe tra i propri 4 migliori. Grazie
alla selezione casuale, prima o poi un peer (Bob) le invia qualche chunk; Alice diventa così capace di
ricambiare, magari entra tra i migliori di Bob, e Bob tra i suoi. Chi carica di più riceve di più: i
``parassiti'' che scaricano senza caricare vengono penalizzati.
\end{esempio}

% ══════════════════════════════════════════════════════════════
\section{Riepilogo}
% ══════════════════════════════════════════════════════════════

\riepilogo{
\begin{itemize}
    \item La posta usa \textbf{user agent} e \textbf{mail server}; \textbf{SMTP} (TCP, porta 25) trasferisce i
          messaggi fino alla casella del destinatario, anche tra server.
    \item Per leggere la posta servono protocolli di accesso: \textbf{POP3} (semplice: autorizzazione,
          transazione, aggiornamento), \textbf{IMAP} (cartelle e ricerca sul server), \textbf{HTTP} (webmail).
    \item Nel \textbf{P2P} i peer sono client e server insieme: la distribuzione di file scala meglio.
    \item \textbf{BitTorrent}: torrent, chunk, \textbf{tracker}; \textbf{rarest first} per scaricare,
          \textbf{4 migliori ogni 10 s} + \textbf{uno casuale ogni 30 s} per caricare.
\end{itemize}
}
